\section{Ruta Planificada: utilidad dinámica por ruta planificada}
\label{sec:exp_e2}
%%%%%%%%%%%%%%%%%%%%%%%%%%%%%%%%%%%%%%%%%%%%%%%%%%%%%%%%%%%%%%%%%%%%%%%%%%%%%%%%

\textcolor{red}{%
[SCAFFOLDING --- Esta sección presenta los experimentos de validación para la estrategia Ruta Planificada, que calcula el costo dinámico $D(v)$ contando el número de personas detectadas dentro de un radio $r$ de cualquier punto de la ruta planificada hacia la frontera $v$:
\[
    D_{\text{Ruta}}(v) = \bigl|\{ p_i : \exists\, q \in \text{ruta}(v),\; \| q - p_i \|_2 \leq r \}\bigr|
\]
A diferencia de la estrategia Vecindad, la estrategia Ruta Planificada evalúa la interferencia a lo largo de todo el camino que recorrerá el robot, en lugar de solo en la vecindad del destino final.
Los casos de prueba son TC011–TC020, con tres réplicas por caso (V031–V060).
Los hallazgos principales son: Ruta Planificada supera a Greedy en celdas fantasma en 5/10 configuraciones, en entropía en 4/10 y en frames en 3/10; por puntuación compuesta, Ruta Planificada es competitiva o mejor en 5/10 configuraciones.
--- Completar con análisis detallado, tablas y figuras de Ruta Planificada vs Greedy.%
]}

%%%%%%%%%%%%%%%%%%%%%%%%%%%%%%%%%%%%%%%%%%%%%%%%%%%%%%%%%%%%%%%%%%%%%%%%%%%%%%%%
